1.1. UV continuum slope ($\beta$) \hfill(23 marks)

\begin{description}
\item[(D1.1.1)] \hfill(6 marks)\\
  The AB magnitude of a typical galaxy is roughly constant, therefore:
  \[ f_\nu = \mathit{const.} \quad\ldots\ldots(1) \]
  Conversion between $f_\nu$ and $f_\lambda$ could be expressed as:
  \[ f_\nu\,\Delta\nu = f_\nu\left(\frac{c}{\lambda-\Delta\lambda}
     - \frac{c}{\lambda}\right)
     = f_\nu\,\frac{c}{\lambda^2}\,\Delta\lambda = f_\lambda\,\Delta\lambda
     \;\Rightarrow\; f_\lambda \propto \lambda^{-2} \quad\ldots\ldots(2) \]
  Therefore, the UV continuum slope should be $\beta = -2$.

\item[(D1.1.2)] \hfill(12 marks)\\
  First of all, we need to convert all the wavelength into rest-frame value,
  using the definition of redshift:
  \[ z = \frac{\lambda_{\mathrm{obs}}}{\lambda_{\mathrm{rest}}} - 1
     \;\Rightarrow\;
     \lambda_0 = \frac{\lambda_{\mathrm{obs}}}{1+z} \quad\ldots\ldots(3) \]
  Thus, the rest-frame wavelengths of the four bands are:

  \begin{center}
  \begin{tabular}{lcccc}
    Band & Y & J & H & K \\
    \hline
    Rest-frame Wavelength (\AA) & 1382 & 1645 & 2171 & 2829 \\
  \end{tabular}
  \end{center}

  The plot is shown as Figure 1.
  % FIGURE NEEDED: Figure 1 -- AB magnitude of CR7 vs rest-frame wavelength
  % (log scale) with error bars and best-fit UV continuum line, annotated
  % beta = -2.7, m_1600 = 24.75 (2018_18_da_sol.pdf, page 4)

\item[(D1.1.3)] \hfill(5 marks)\\
  Here we can express the AB magnitude of CR7 (UV continuum only) using
  $\beta$ and $\lambda_{\mathrm{rest}}$:
  \[ m_{AB} = -2.5\log f_\nu + C = -2.5\log(\lambda^2 f_\lambda) + C
     = -2.5(\beta+2)\log\!\left(\frac{\lambda}{1600\,\text{\AA}}\right)
       + m_{1600} \quad\ldots\ldots(4) \]
  Here $m_{1600}$ is the AB magnitude at rest-frame 1600\,\AA\
  (observed-frame $(1+z)\cdot 1600$\,\AA). After writing this out, we can fit
  the data with a linear function $f(\log\lambda) = a\log\lambda + b$, where
  the slope $a = -2.5(\beta+2)$. Least squares fitting result is $a = 1.7$,
  therefore $\beta = -2.7$.

  The best-fit UV continuum should be plot as a straight line on Figure 1.

  Compared to the galaxy in (D1.1), CR7 shows smaller $\beta$ value,
  indicating its UV continuum is bluer than typical galaxy. Therefore, there
  should be less dust content in CR7, and the correct answer should be [NO].
  \hfill(2 marks)
\end{description}

1.2. UV slope and IRX \hfill(20 marks)

\begin{description}
\item[(D1.2.1)] \hfill(14 marks)\\
  For the galaxy sample in Table 2, the IRX could be calculated as:
  \[ \mathrm{IRX} = \log F_{FIR} - \log F_{1600} \quad\ldots\ldots(5) \]
  Therefore, the final dataset we need for plotting is:

  \begin{center}
  \begin{tabular}{lrr|lrr}
    Galaxy Name & $\beta$ & IRX & Galaxy Name & $\beta$ & IRX \\
    \hline
    NGC 4861     & $-2.46$ & $-0.08$ & NGC 7250 & $-1.33$ & 0.46 \\
    Mrk 153      & $-2.41$ & $-0.55$ & NGC 7714 & $-1.23$ & 0.84 \\
    Tol 1924-416 & $-2.12$ & $-0.12$ & NGC 3049 & $-1.14$ & 0.85 \\
    UGC 9560     & $-2.02$ & $-0.03$ & NGC 3310 & $-1.05$ & 1.01 \\
    NGC 3991     & $-1.91$ & 0.34    & NGC 2782 & $-0.9$  & 1.17 \\
    Mrk 357      & $-1.8$  & 0.21    & NGC 1614 & $-0.76$ & 2.07 \\
    Mrk 36       & $-1.72$ & $-0.26$ & NGC 6052 & $-0.72$ & 1.14 \\
    NGC 4670     & $-1.65$ & 0.17    & NGC 3504 & $-0.56$ & 1.45 \\
    NGC 3125     & $-1.49$ & 0.55    & NGC 4194 & $-0.26$ & 1.63 \\
    UGC 3838     & $-1.41$ & 0.26    & NGC 3256 & 0.16    & 1.88 \\
  \end{tabular}
  \end{center}

  The plot is shown as Figure 2. Best-fit function is:
  \[ \mathrm{IRX} = 0.98\cdot\beta + 1.96 \quad\ldots\ldots(6) \]
  which should be plotted as a straight line on Figure 2.
  % FIGURE NEEDED: Figure 2 -- IRX vs beta diagram of the 20 local galaxies
  % with best-fit line IRX = 0.98 beta + 1.96 (2018_18_da_sol.pdf, page 5)

\item[(D1.2.2)] \hfill(6 marks)\\
  With the IRX--$\beta$ relation we derived above, we can calculate the
  difference between observed IRX and expected value (inferred by $\beta$) of
  each galaxy:

  \begin{center}
  \begin{tabular}{lrrr|lrrr}
    Name & $\mathrm{IRX}_{\mathrm{obs}}$ & $\mathrm{IRX}_{\mathrm{exp}}$ &
    $\Delta$IRX & Name & $\mathrm{IRX}_{\mathrm{obs}}$ &
    $\mathrm{IRX}_{\mathrm{exp}}$ & $\Delta$IRX \\
    \hline
    NGC 4861     & $-0.08$ & $-0.45$ & 0.37    & NGC 7250 & 0.46 & 0.66 & $-0.20$ \\
    Mrk 153      & $-0.55$ & $-0.40$ & $-0.15$ & NGC 7714 & 0.84 & 0.75 & 0.09 \\
    Tol 1924-416 & $-0.12$ & $-0.12$ & 0.00    & NGC 3049 & 0.85 & 0.84 & 0.01 \\
    UGC 9560     & $-0.03$ & $-0.02$ & $-0.01$ & NGC 3310 & 1.01 & 0.93 & 0.08 \\
    NGC 3991     & 0.34    & 0.09    & 0.25    & NGC 2782 & 1.17 & 1.08 & 0.09 \\
    Mrk 357      & 0.21    & 0.20    & 0.01    & NGC 1614 & 2.07 & 1.22 & 0.85 \\
    Mrk 36       & $-0.26$ & 0.27    & $-0.53$ & NGC 6052 & 1.14 & 1.25 & $-0.11$ \\
    NGC 4670     & 0.17    & 0.34    & $-0.17$ & NGC 3504 & 1.45 & 1.41 & 0.04 \\
    NGC 3125     & 0.55    & 0.50    & 0.05    & NGC 4194 & 1.63 & 1.71 & $-0.08$ \\
    UGC 3838     & 0.26    & 0.58    & $-0.32$ & NGC 3256 & 1.88 & 2.12 & $-0.24$ \\
  \end{tabular}
  \end{center}

  Therefore, the dispersion of IRX is:
  \[ \sigma = \sqrt{\frac{\sum_{i=1}^{N}(\Delta\mathrm{IRX}_i)^2}{N-1}}
     = 0.28\ \mathrm{dex} \quad\ldots\ldots(7) \]
\end{description}

1.3. IRX and $A_{1600}$ \hfill(20 marks)

\begin{description}
\item[(D1.3.1)] \hfill(6 marks)\\
  Based on the two assumptions we made in Question (3), we can simplify
  $F_{FIR}/F_{1600}$ as:
  \[ \frac{F_{FIR}}{F_{1600}} \approx 10^{0.4 A_{1600}} - 1
     \quad\ldots\ldots(8) \]
  \[ \therefore A_{1600} = 2.5\log\left(1 + 10^{\mathrm{IRX}}\right)
     \quad\ldots\ldots(9) \]

\item[(D1.3.2)] \hfill(12 marks)\\
  Now we can calculate $A_{1600}$ of all the galaxies in Table 2 with
  Equation 9:

  \begin{center}
  \begin{tabular}{lrr|lrr}
    Name & $\beta$ & $A_{1600}$ & Name & $\beta$ & $A_{1600}$ \\
    \hline
    NGC 4861     & $-2.46$ & 0.66 & NGC 7250 & $-1.33$ & 1.47 \\
    Mrk 153      & $-2.41$ & 0.27 & NGC 7714 & $-1.23$ & 2.25 \\
    Tol 1924-416 & $-2.12$ & 0.61 & NGC 3049 & $-1.14$ & 2.27 \\
    UGC 9560     & $-2.02$ & 0.72 & NGC 3310 & $-1.05$ & 2.63 \\
    NGC 3991     & $-1.91$ & 1.26 & NGC 2782 & $-0.9$  & 3.00 \\
    Mrk 357      & $-1.8$  & 1.05 & NGC 1614 & $-0.76$ & 5.18 \\
    Mrk 36       & $-1.72$ & 0.48 & NGC 6052 & $-0.72$ & 2.93 \\
    NGC 4670     & $-1.65$ & 0.99 & NGC 3504 & $-0.56$ & 3.66 \\
    NGC 3125     & $-1.49$ & 1.64 & NGC 4194 & $-0.26$ & 4.10 \\
    UGC 3838     & $-1.41$ & 1.13 & NGC 3256 & 0.16    & 4.71 \\
  \end{tabular}
  \end{center}

  Then we can plot the $A_{1600}$--$\beta$ relation as Figure 3. The best-fit
  equation is:
  \[ A_{1600} = 1.92\cdot\beta + 4.62 \quad\ldots\ldots(10) \]
  % FIGURE NEEDED: Figure 3 -- A_1600 vs beta diagram with best-fit line
  % A_1600 = 1.92 beta + 4.62 (2018_18_da_sol.pdf, page 7)

\item[(D1.3.3)] \hfill(2 marks)\\
  If Equation 10 is correct, then for a dust-free galaxy, the dust
  attenuation $A_{1600}$ should equal to 0. Therefore, we can derive
  $\beta_0 \approx -2.4$, which is the expect UV continuum slope of a
  ``dust-free'' galaxy.

  (Note: The $\beta_0$ result here is conflict with some observations, for
  example, the UV slope of CR7 which we just derived. In fact, the linear
  model here is not necessarily correct and it is just a simplification of
  complex model: we have made two very rough assumption to derive the
  $A_{1600}-\mathrm{IRX}$ relation, and other physical parameters of galaxy,
  e.g.\ metallicity, will also influence the overall $A_{1600}-\beta$
  relation.)
\end{description}

1.4. Extend to high-redshift universe \hfill(12 marks) (discarded)

\begin{description}
\item[(D1.4.1)] \hfill(6 marks)\\
  Since we have already known the upper limit of $F_{FIR}$ of CR7, we only
  need to derive $F_{1600}$ in order to calculate IRX. First of all, we need
  to translate the $m_{1600}$ we got in (D 1.3) to
  $f_{(1+z)1600\text{\AA}}$:
  \[ f_{\lambda_0(1+z)} = f_{\nu_0(1+z)}\cdot
     \frac{c}{[\lambda_0(1+z)]^2}
     = \frac{c}{[\lambda_0(1+z)]^2}\cdot 10^{-0.4 m_{1600}}\cdot
     3631\ \mathrm{Jy} \quad\ldots\ldots(11) \]
  and we will derive
  $f_{\lambda_0(1+z)} = 9.3\times 10^{-13}\,
  \mathrm{J\,s^{-1}\,m^{-3}}$. Therefore:
  \[ F_{1600} = (1+z)\cdot\lambda_0\cdot f_{\lambda_0\cdot(1+z)}
     = 1.1\times 10^{-18}\,\mathrm{J\,s^{-1}\,m^{-2}}
     \quad\ldots\ldots(12) \]
  Thus, the upper limit of CR7's IRX should be:
  \[ \mathrm{IRX}_{CR7} = \log\frac{F_{FIR}}{F_{1600}} = -0.87
     \quad\ldots\ldots(13) \]
  It is an upper limit.

\item[(D1.4.2)] \hfill(6 marks)\\
  Given the UV continuum slop of CR7 ($-2.7$), we can infer its expected IRX
  from the IRX$\,-\,\beta$ relation in local universe:
  \[ \widehat{\mathrm{IRX}}_{CR7} = 0.98\cdot\beta_{CR7} + 1.96 = -0.69
     \quad\ldots\ldots(14) \]
  The upper limit of observed IRX is 0.18\,dex lower than the value predicted
  by local relation. However, we have already known the local relation has a
  dispersion of 0.28\,dex, which is larger than the 0.18\,dex difference.
  This indicates that the current IRX upper limit is not tight enough (in
  other words, current observation is not deep enough) to show any difference
  of CR7 from local galaxies on the diagram of IRX$-\beta$. Therefore, the
  correct answer should be [NO].
\end{description}
